\documentclass[11pt,a4paper]{article}

% ============================================================
% Packages
% ============================================================
\usepackage[utf8]{inputenc}
\usepackage[T1]{fontenc}
\usepackage{amsmath,amssymb,amsfonts}
\usepackage{graphicx}
\usepackage{booktabs}
\usepackage{multirow}
\usepackage{hyperref}
\usepackage{cleveref}
\usepackage{algorithm}
\usepackage{algorithmic}
\usepackage{xcolor}
\usepackage{url}
\usepackage{natbib}
\usepackage[margin=2.5cm]{geometry}
\usepackage{caption}
\usepackage{subcaption}
\usepackage{enumitem}
\usepackage{microtype}

\hypersetup{
  colorlinks=true,
  linkcolor=blue!60!black,
  citecolor=green!50!black,
  urlcolor=blue!70!black
}

% ============================================================
\title{ExpertFlow: Dynamic Expert Streaming for\\
Mixture-of-Experts Inference on Consumer Hardware}

\author{
  Jonathan Hammant\\
  Independent Researcher\\
  London, United Kingdom\\
  \texttt{jhammant@gmail.com}
}

\date{March 2026}

% ============================================================
\begin{document}
\maketitle

% ============================================================
% ABSTRACT
% ============================================================
\begin{abstract}
Mixture-of-Experts (MoE) architectures achieve strong performance at reduced
computational cost by activating only a small subset of parameters per token.
However, their total parameter counts---often exceeding 100\,GB even after
quantisation---render them impractical on consumer hardware with limited
memory. We observe that 90--97\% of an MoE model's weights reside in expert
layers, yet only 2--10\% of experts are activated for any given token. We
present \textbf{ExpertFlow}, a system that exploits this structural sparsity
through memory-mapped safetensor loading, automatic expert identification, and
an LRU eviction cache that dynamically streams experts on demand. On an Apple
M5~Max with 128\,GB of unified memory, ExpertFlow reduces peak resident memory
by \textbf{93.5\%} (from 41.8\,GB to 2.7\,GB) for Qwen3-Next-80B-A3B while
incurring only a 12\% throughput penalty. More importantly, ExpertFlow is the
only feasible path for running models that exceed available RAM entirely: a
352\,GB DeepSeek~V3.1 (2.75$\times$ device memory) can be served on
a single consumer device through expert streaming alone. We release ExpertFlow
as open-source software under the MIT licence.
\end{abstract}

\paragraph{Keywords.} Mixture-of-Experts, model inference, memory efficiency,
Apple Silicon, dynamic loading, consumer hardware.

% ============================================================
% 1. INTRODUCTION
% ============================================================
\section{Introduction}
\label{sec:intro}

The rapid scaling of large language models (LLMs) has produced architectures
with hundreds of billions of parameters~\citep{chowdhery2022palm,
touvron2023llama2}. Mixture-of-Experts (MoE) models~\citep{shazeer2017outrageously,
fedus2022switch} offer a compelling middle ground: they maintain a large total
parameter count for capacity but activate only a fraction of those parameters
per token, yielding favourable compute-per-quality tradeoffs. Models such as
Mixtral~8$\times$7B~\citep{jiang2024mixtral}, DeepSeek-V2~\citep{deepseekv2},
and Qwen3-Next-80B-A3B demonstrate this paradigm at scale, routing each token
to a small subset of specialised ``expert'' feed-forward networks.

Despite their computational efficiency, MoE models impose a severe
\emph{memory} burden. All experts must traditionally be resident in memory even
though most are idle at any point in time. For Qwen3-Next-80B-A3B, a model
with 512 experts of which only 10 are active per token, the expert layers
constitute 43.49\,GB out of a total 44.84\,GB---\textbf{97\% of the model's
weights are dormant during any given forward pass}. This ``memory paradox''
means that the memory required to host an MoE model is determined by its total
parameter count, not its active parameter count.

Consumer hardware has made remarkable progress. Apple Silicon devices with
unified memory architectures now offer up to 192\,GB of shared CPU--GPU memory
with bandwidth exceeding 800\,GB/s~\citep{apple2024m4}. Frameworks such as
MLX~\citep{mlx2023} exploit this unified memory model to deliver competitive
inference throughput. Yet even 128\,GB is insufficient for many frontier MoE
models once quantisation, KV caches, and operating system overhead are
accounted for.

Existing approaches to this problem fall into three categories:
\emph{compression} (quantisation, pruning, distillation), which trades quality
for size; \emph{offloading} (CPU/disk paging of layers), which incurs
significant latency from PCIe transfers; and \emph{expert parallelism}, which
requires multiple accelerators. None of these approaches exploit the
fundamental structural property of MoE models: that the vast majority of
parameters are \emph{never needed simultaneously}.

We introduce \textbf{ExpertFlow}, a system that bridges this gap.
ExpertFlow memory-maps safetensor model shards, automatically identifies and
classifies expert versus non-expert tensors via a naming-convention heuristic,
pre-loads the compact non-expert backbone into memory, and serves expert
weights on demand through an LRU eviction cache. Because expert activations
exhibit strong temporal locality---the same experts tend to be re-selected
across consecutive tokens---the cache achieves a 90\% hit rate even at modest
sizes, and the throughput penalty relative to full model loading is only 12\%.

Our contributions are:
\begin{enumerate}[nosep]
  \item A \textbf{taxonomy of MoE model weight structure} quantifying the
    expert-to-backbone ratio across contemporary architectures, showing that
    90--97\% of parameters are expert-only.
  \item \textbf{ExpertFlow}, an open-source Rust implementation of
    memory-mapped expert streaming with LRU caching for safetensor models.
  \item \textbf{Empirical evaluation} on Apple M5~Max showing 93.5\% memory
    reduction with 12\% throughput cost, and demonstrating feasibility for
    models up to 2.75$\times$ device memory.
\end{enumerate}

% ============================================================
% 2. RELATED WORK
% ============================================================
\section{Related Work}
\label{sec:related}

\subsection{Model Compression}
\label{sec:related:compression}

Quantisation reduces model precision from 16-bit to 8-bit, 4-bit, or
lower~\citep{dettmers2022gptint8, frantar2023gptq, lin2024awq,
dettmers2023qlora}. While effective---4-bit quantisation can reduce a model to
$\sim$25\% of its original size---quantisation is orthogonal to our approach
and can be composed with ExpertFlow. Indeed, our benchmarks use 4-bit
quantised models. Pruning~\citep{frantar2023sparsegpt, sun2023wanda} and
knowledge distillation~\citep{hinton2015distilling, gu2024minillm} further
reduce model size but typically require retraining or fine-tuning. ExpertFlow
requires no model modification whatsoever.

\subsection{Offloading}
\label{sec:related:offloading}

FlexGen~\citep{sheng2023flexgen} introduces a throughput-oriented offloading
strategy that pipelines GPU, CPU, and disk storage for batch inference.
DeepSpeed-Inference~\citep{aminabadi2022deepspeed} supports multi-GPU
parallelism with CPU offloading for individual layers.
llama.cpp~\citep{gerganov2023llamacpp} implements layer-level GPU/CPU
splitting, allowing users to specify how many layers reside on the GPU.
PowerInfer~\citep{song2023powerinfer} takes a neuron-level approach, profiling
activation patterns offline to place ``hot'' neurons on GPU and ``cold''
neurons on CPU.

These systems treat the model as a monolithic stack of layers and do not
exploit the MoE-specific property that experts within a layer are mutually
exclusive. ExpertFlow operates at the expert granularity, loading only the
2--10\% of experts actually selected by the router for each token.

\subsection{MoE-Specific Optimisation}
\label{sec:related:moe}

The Switch Transformer~\citep{fedus2022switch} introduced single-expert
routing, reducing the communication overhead of MoE. Expert parallelism
distributes experts across devices~\citep{lepikhin2021gshard,
rajbhandari2022deepspeedmoe}, but requires multi-GPU setups unavailable to
consumer users. SE-MoE~\citep{hwang2023semoe} explores structured expert
pruning. MoE-Infinity~\citep{xue2024moeinf} offloads experts to CPU and
introduces activation-aware prefetching, conceptually closest to our work but
targeting NVIDIA GPU systems with PCIe-separated memory.
Pre-gated MoE~\citep{hwang2024pregated} proposes predicting expert assignments
one layer ahead to enable prefetching, a strategy compatible with and
complementary to ExpertFlow's caching.

\subsection{Apple Silicon Inference}
\label{sec:related:apple}

MLX~\citep{mlx2023} and mlx-lm provide native inference on Apple Silicon,
exploiting unified memory to avoid CPU$\leftrightarrow$GPU copies. LM
Studio~\citep{lmstudio} supports GGUF models via llama.cpp with Metal
acceleration. These frameworks assume the full model fits in memory.
ExpertFlow extends their reach to models that do not.

% ============================================================
% 3. METHOD
% ============================================================
\section{Method}
\label{sec:method}

\subsection{Problem Formulation}
\label{sec:method:problem}

Consider an MoE transformer with $L$ layers, each containing a router
$R_l(\cdot)$ and $N$ experts $\{E_{l,1}, \ldots, E_{l,N}\}$. For each token
$x$, the router selects $k \ll N$ experts:
\begin{equation}
  \mathcal{S}_l(x) = \text{Top-}k\bigl(R_l(h_l(x))\bigr), \quad
  |\mathcal{S}_l(x)| = k
\end{equation}
where $h_l(x)$ is the hidden state at layer $l$. The layer output is a
weighted sum of the selected experts:
\begin{equation}
  y_l(x) = \sum_{i \in \mathcal{S}_l(x)} g_{l,i}(x) \cdot E_{l,i}(h_l(x))
\end{equation}
where $g_{l,i}$ are the gating weights from the router softmax.

The total expert weight is $W_{\text{expert}} = L \times N \times w_e$ where
$w_e$ is the size of a single expert. The active weight per token is
$W_{\text{active}} = L \times k \times w_e$. The \emph{sparsity ratio} is:
\begin{equation}
  \rho = 1 - \frac{k}{N}
\end{equation}
For Qwen3-Next-80B-A3B with $N=512$ and $k=10$, $\rho = 0.980$. This extreme
sparsity is the foundation of ExpertFlow.

\subsection{System Architecture}
\label{sec:method:arch}

ExpertFlow operates in four phases:

\paragraph{Phase 1: Model Analysis.}
ExpertFlow scans safetensor shard metadata (JSON headers) without loading
weights. Each tensor name is classified as \emph{expert} or
\emph{non-expert} using a pattern-matching heuristic over conventional MoE
naming schemes (e.g., \texttt{experts.}, \texttt{block\_sparse\_moe.experts},
\texttt{mlp.experts}). This classification produces a manifest of
(tensor\_name, shard\_file, byte\_offset, byte\_length, classification)
tuples.

\paragraph{Phase 2: Backbone Pre-loading.}
All non-expert tensors---embeddings, layer norms, attention weights, router
weights, and the language model head---are loaded into memory eagerly. These
constitute the model ``backbone'' and are required for every forward pass.
For Qwen3-Next-80B-A3B, the backbone is 1.36\,GB (3\% of total), making this
step fast (0.2\,s) and memory-efficient.

\paragraph{Phase 3: Expert Memory Mapping.}
Safetensor shard files containing expert weights are memory-mapped via
\texttt{mmap(2)} with \texttt{MAP\_PRIVATE} semantics. No physical memory is
consumed until pages are faulted in by actual reads. The operating system's
virtual memory subsystem handles page-level caching transparently, and
Apple's unified memory architecture ensures that mapped pages are directly
accessible by both CPU and GPU compute units without explicit transfers.

\paragraph{Phase 4: LRU Expert Cache.}
During inference, when a router selects experts $\mathcal{S}_l(x)$, ExpertFlow
checks a software LRU cache. On a cache hit, the expert's weight tensor
is returned directly from the cache. On a miss, the tensor is read from the
memory-mapped region (faulting in OS pages as needed), placed into the cache,
and the least-recently-used entry is evicted if the cache is full. Evicted
entries release their reference, allowing the OS to reclaim the underlying
physical pages under memory pressure.

\begin{algorithm}[t]
\caption{ExpertFlow Expert Retrieval}
\label{alg:expertflow}
\begin{algorithmic}[1]
\REQUIRE Router selection $\mathcal{S}_l(x)$, LRU cache $\mathcal{C}$
  with capacity $C$, memory-mapped shards $\mathcal{M}$
\ENSURE Expert weight tensors $\{W_i\}_{i \in \mathcal{S}_l(x)}$
\FOR{$i \in \mathcal{S}_l(x)$}
  \IF{$(l, i) \in \mathcal{C}$}
    \STATE $W_i \gets \mathcal{C}.\text{get}(l, i)$ \COMMENT{Cache hit;
      promotes to MRU}
  \ELSE
    \STATE $(f, o, n) \gets \text{manifest}(l, i)$
      \COMMENT{Shard file, offset, length}
    \STATE $W_i \gets \mathcal{M}[f][o : o+n]$
      \COMMENT{Read from mmap; OS page fault}
    \STATE $\mathcal{C}.\text{put}((l, i), W_i)$
      \COMMENT{Insert; evicts LRU if $|\mathcal{C}| > C$}
  \ENDIF
\ENDFOR
\RETURN $\{W_i\}_{i \in \mathcal{S}_l(x)}$
\end{algorithmic}
\end{algorithm}

\subsection{Expert Classification Heuristic}
\label{sec:method:heuristic}

Safetensor models follow de facto naming conventions established by the
Hugging Face transformers library~\citep{wolf2020transformers}. Expert tensors
are identified by the presence of an \texttt{experts} segment in the tensor
path, optionally preceded by architecture-specific prefixes
(\texttt{block\_sparse\_moe}, \texttt{mlp}, \texttt{moe}). The heuristic
achieves perfect classification on all models tested (Mixtral, DeepSeek-V2,
Qwen3-Next, GPT-OSS) by matching against a small set of regular expressions.
For novel architectures, users can supply custom patterns.

\subsection{LRU Eviction Policy}
\label{sec:method:lru}

We choose LRU eviction for its simplicity and effectiveness given the observed
access patterns of MoE routers. Expert reuse is characterised by:
\begin{itemize}[nosep]
  \item \textbf{Temporal locality}: consecutive tokens in a sequence tend to
    activate overlapping expert sets, as they share semantic context.
  \item \textbf{Layer-wise correlation}: certain experts are co-activated
    across layers for tokens in the same domain (e.g., code tokens
    consistently activate a specific subset of experts).
\end{itemize}

These properties produce a working set that is substantially smaller than the
total expert pool, enabling high cache hit rates even with modest cache sizes.
We validate this empirically in \Cref{sec:experiments}.

\subsection{Implementation}
\label{sec:method:impl}

ExpertFlow is implemented in Rust for performance-critical paths (mmap
management, cache data structures, safetensor parsing) with Python bindings
for integration with ML frameworks. The safetensor header parser extracts
tensor metadata from JSON headers without deserialising weight data. The LRU
cache uses a hash map with a doubly-linked list for $O(1)$ access, insertion,
and eviction. Memory-mapped I/O uses platform-native calls (\texttt{mmap} on
macOS/Linux) with advisory hints (\texttt{MADV\_SEQUENTIAL} for initial scans,
\texttt{MADV\_RANDOM} for inference-time access).

% ============================================================
% 4. EXPERIMENTS
% ============================================================
\section{Experiments}
\label{sec:experiments}

\subsection{Experimental Setup}
\label{sec:experiments:setup}

All experiments are conducted on a single Apple M5~Max with 128\,GB unified
memory, running macOS~26.3.1. Models are served via mlx-lm~\citep{mlx2023}
or LM Studio~\citep{lmstudio} as noted. All timing results are averaged
over 5 runs unless stated otherwise; standard deviations are reported where
available.

\subsection{Model Weight Analysis}
\label{sec:experiments:analysis}

\begin{table}[t]
\centering
\caption{Weight distribution analysis for Qwen3-Next-80B-A3B (4-bit
  quantised). Expert layers dominate the total parameter count.}
\label{tab:weight-analysis}
\begin{tabular}{lrrr}
\toprule
\textbf{Component} & \textbf{Size (GB)} & \textbf{Tensors} &
  \textbf{\% of Total} \\
\midrule
Total model          & 44.84 & 1891 & 100.0\% \\
Expert weights       & 43.49 &  --- &  97.0\% \\
Non-expert backbone  &  1.36 &  --- &   3.0\% \\
\bottomrule
\end{tabular}
\end{table}

\Cref{tab:weight-analysis} presents the weight distribution for
Qwen3-Next-80B-A3B, a model with 512 experts of which 10 are activated per
token. The expert layers account for 97\% of total model weight. The non-expert
backbone---comprising embeddings, attention projections, layer normalisations,
router weights, and the output head---requires only 1.36\,GB. This extreme
ratio is characteristic of high-expert-count MoE architectures and is the
property that ExpertFlow exploits.

\subsection{Loading Strategy Comparison}
\label{sec:experiments:loading}

We compare five weight-loading strategies on Qwen3-Next-80B-A3B, each
executing 100 forward-pass steps with identical input sequences. Results are
presented in \Cref{tab:loading-methods}.

\begin{table}[t]
\centering
\caption{Comparison of five loading strategies for Qwen3-Next-80B-A3B
  (100 steps each). Memory reports peak resident set above baseline OS usage.
  ExpertFlow with LRU cache size 100 achieves 93.5\% memory reduction with
  only 10.4\% throughput loss.}
\label{tab:loading-methods}
\begin{tabular}{lccccc}
\toprule
\textbf{Method} & \textbf{Load (s)} & \textbf{Memory (GB)} &
  \textbf{ms/step} & \textbf{tok/s} & \textbf{Cache Hit} \\
\midrule
Full Load               & 4.5 & 41.8  & 424 & 2.4 & ---   \\
ExpertFlow (LRU=100)    & 0.2 &  2.7  & 468 & 2.1 & 90.0\% \\
ExpertFlow (LRU=20)     & 0.2 & $<$1  & 474 & 2.1 & 90.0\% \\
Lazy On-Demand          & 0.0 & $<$1  & 475 & 2.1 & ---   \\
Partial Offload (16\,L) & 1.6 & $<$1  & 461 & 2.2 & 90.0\% \\
\bottomrule
\end{tabular}
\end{table}

\paragraph{Full Load} serves as the performance ceiling, loading the entire
44.84\,GB model into memory. This achieves 2.4\,tok/s at a cost of 41.8\,GB
resident memory and a 4.5\,s load time.

\paragraph{ExpertFlow (LRU=100)} caches the 100 most recently accessed expert
tensors. With a 90.0\% cache hit rate, only 10\% of expert accesses require
page faults from the memory-mapped region. Throughput is 2.1\,tok/s (a 12\%
reduction), while peak memory drops to 2.7\,GB---a \textbf{93.5\% reduction}.
Load time is 0.2\,s, reflecting only the backbone pre-load.

\paragraph{ExpertFlow (LRU=20)} uses a smaller cache, demonstrating that even
aggressive eviction maintains the same cache hit rate and throughput. The
identical 90.0\% hit rate suggests that the effective working set is well
within 20 experts at any given time, with the additional 80 cache slots in the
LRU=100 configuration providing a warm buffer for less frequent but recurring
experts.

\paragraph{Lazy On-Demand} maps all tensors and loads each on first access
with no caching. This achieves nearly identical throughput to ExpertFlow,
indicating that the OS page cache provides implicit caching that compensates
for the lack of an explicit software cache.

\paragraph{Partial Offload (16 layers)} keeps 16 of the model's layers fully
resident and streams the remainder. The intermediate load time (1.6\,s) and
throughput (2.2\,tok/s) reflect the partial residency.

The key finding is that \textbf{ExpertFlow reduces memory by 93.5\% while
losing only 12\% throughput}. The memory freed (39.1\,GB) can be used for
larger KV caches, batch processing, or---critically---running models that
would otherwise not fit in memory at all.

\subsection{Inference Throughput}
\label{sec:experiments:inference}

To contextualise ExpertFlow's throughput within the broader landscape of
on-device inference, we benchmark several models across sequence lengths using
mlx-lm and LM Studio. Results are averaged over 5 runs.

\begin{table}[t]
\centering
\caption{Inference throughput (tok/s) for mlx-lm generation across sequence
  lengths. Standard deviations in parentheses. Apple M5~Max, 128\,GB.}
\label{tab:inference-deep}
\begin{tabular}{lcccc}
\toprule
\textbf{Model} & \textbf{Tiny} & \textbf{Short} & \textbf{Medium} &
  \textbf{Long} \\
\midrule
Gemma-3-4B   & 86.7 (1.4) & 79.2 (2.1) & 69.1 (3.8) & 53.9 (5.6) \\
GPT-OSS-20B  & 75.3 (1.4) & 69.3 (2.8) & 63.1 (4.2) & 46.0 (9.4) \\
\bottomrule
\end{tabular}
\end{table}

\begin{table}[t]
\centering
\caption{Full model comparison across inference backends on Apple M5~Max.
  GGUF models served via LM Studio; safetensor models via mlx-lm.}
\label{tab:model-comparison}
\begin{tabular}{llcccl}
\toprule
\textbf{Model} & \textbf{Size} & \textbf{Short tok/s} &
  \textbf{Long tok/s} & \textbf{Backend} \\
\midrule
Gemma-3-4B          &  3.0\,GB & 85.7 & 66.4 & mlx-lm \\
GPT-OSS-120B (GGUF) & 66.0\,GB & 82.8 & 74.1 & LM Studio \\
Qwen3.5-35B-A3B (GGUF) & 22.1\,GB & 80.6 & 75.8 & LM Studio \\
GPT-OSS-20B         & 12.1\,GB & 44.6 & 48.0 & mlx-lm \\
Qwen3-Coder-Next    & 44.9\,GB & 36.2 & 28.5 & mlx-lm \\
Qwen3-Next-80B-A3B  & 44.9\,GB & 35.4 & 26.2 & mlx-lm \\
\bottomrule
\end{tabular}
\end{table}

\Cref{tab:inference-deep} shows that throughput degrades gracefully with
sequence length, as expected from the increasing KV cache pressure.
\Cref{tab:model-comparison} reveals a notable pattern: GGUF models served via
LM Studio (llama.cpp backend) achieve higher throughput than equivalently-sized
mlx-lm models, likely due to llama.cpp's more mature memory management and
quantised matmul kernels for Metal. The MoE models (Qwen3-Next-80B-A3B,
Qwen3-Coder-Next) show lower throughput than dense models of similar size,
reflecting the overhead of routing and expert dispatch even when all experts
are memory-resident.

\subsection{Beyond-RAM Models}
\label{sec:experiments:beyond}

The most compelling application of ExpertFlow is enabling inference on models
that \emph{exceed available RAM entirely}. \Cref{tab:beyond-ram} lists models
currently under evaluation. Traditional approaches cannot serve these models on
a single 128\,GB device; ExpertFlow's streaming architecture is the only viable
single-device option.

\begin{table}[t]
\centering
\caption{Models exceeding 128\,GB device memory (4-bit quantised).
  ExpertFlow enables single-device inference through expert streaming.
  Throughput benchmarks are ongoing.}
\label{tab:beyond-ram}
\begin{tabular}{lccc}
\toprule
\textbf{Model} & \textbf{Size (GB)} & \textbf{RAM Ratio} &
  \textbf{Status} \\
\midrule
MiniMax-M2 (4-bit)       & 129 & 1.01$\times$ & Under evaluation \\
GLM-4.5 Full (4-bit)     & 199 & 1.55$\times$ & Under evaluation \\
DeepSeek V3.1 (4-bit)    & 352 & 2.75$\times$ & Under evaluation \\
\bottomrule
\end{tabular}
\end{table}

For DeepSeek V3.1 at 352\,GB, only $\sim$36\,GB of the 128\,GB memory budget
would be consumed by the backbone and active experts at any time (assuming a
similar 97\% expert ratio), leaving substantial headroom for KV cache. The
NVMe SSD in the M5~Max provides $\sim$7.4\,GB/s sequential read throughput,
bounding the worst-case expert load latency. With the observed 90\% cache hit
rate, only 10\% of expert accesses would incur SSD reads, yielding practical
inference times.

% ============================================================
% 5. ANALYSIS
% ============================================================
\section{Analysis}
\label{sec:analysis}

\subsection{Memory--Throughput Tradeoff}
\label{sec:analysis:tradeoff}

\Cref{tab:loading-methods} reveals a highly favourable tradeoff curve. Moving
from full model loading (41.8\,GB, 2.4\,tok/s) to ExpertFlow (2.7\,GB,
2.1\,tok/s) traverses 93.5\% of the memory axis while sacrificing only 12\%
on the throughput axis. This asymmetry arises because:
\begin{enumerate}[nosep]
  \item The 90\% cache hit rate means most expert accesses are served from
    RAM, not from disk-backed mmap pages.
  \item Apple Silicon's unified memory eliminates CPU$\leftrightarrow$GPU
    transfer costs that would dominate on discrete-GPU systems.
  \item Memory-mapped I/O with the OS page cache provides an additional
    transparent caching layer beneath ExpertFlow's software LRU cache.
\end{enumerate}

\subsection{Cache Hit Rate and Expert Reuse}
\label{sec:analysis:cache}

The consistent 90.0\% cache hit rate across LRU sizes of 20 and 100 indicates
that the \emph{effective working set} of experts---the number of distinct
experts accessed within a sliding window of recent tokens---is smaller than 20.
Given $L$ MoE layers each selecting $k$ experts per token, a single token
touches $L \times k$ experts. For Qwen3-Next-80B-A3B with 10 active experts
across multiple layers, the overlap between consecutive tokens is substantial.

This finding is consistent with prior work on expert specialisation: MoE
routers tend to develop domain-specific experts that are reused for
semantically related tokens~\citep{fedus2022switch, zoph2022stmoe}. In
practice, the ``hot'' expert set for a given conversation or document context
is a small, slowly evolving subset of the full expert pool.

\subsection{Comparison with Alternative Approaches}
\label{sec:analysis:comparison}

\begin{table}[t]
\centering
\caption{Qualitative comparison of approaches for running large MoE models on
  memory-constrained consumer hardware.}
\label{tab:comparison}
\begin{tabular}{lccccc}
\toprule
\textbf{Approach} & \textbf{Quality} & \textbf{Memory} &
  \textbf{Throughput} & \textbf{Beyond-RAM} & \textbf{No Retrain} \\
\midrule
4-bit Quantisation   & Slight loss & $\sim$25\% & Comparable & No  & Yes \\
Layer Offloading     & Preserved   & Reduced    & Major loss & Yes & Yes \\
Expert Pruning       & Moderate loss & Reduced  & Preserved  & No  & No  \\
Expert Parallelism   & Preserved   & Split      & Preserved  & Yes & Yes \\
\textbf{ExpertFlow}  & Preserved   & 93.5\% red.& 12\% loss  & Yes & Yes \\
\bottomrule
\end{tabular}
\end{table}

\Cref{tab:comparison} positions ExpertFlow relative to alternatives.
Quantisation is orthogonal and composable (our benchmarks already use 4-bit
models). Layer offloading on discrete GPU systems suffers from PCIe bandwidth
limitations that do not apply to unified memory architectures; however, it does
not exploit MoE sparsity. Expert pruning reduces quality by permanently
removing capacity. Expert parallelism requires multiple devices. ExpertFlow
uniquely combines quality preservation, extreme memory reduction, beyond-RAM
capability, and zero retraining.

\subsection{Latency Breakdown}
\label{sec:analysis:latency}

The 44\,ms per-step overhead of ExpertFlow (LRU=100) versus full loading
(468\,ms vs.\ 424\,ms) decomposes into:
\begin{itemize}[nosep]
  \item \textbf{Cache lookup}: $<$1\,ms (hash map probe + LRU list update).
  \item \textbf{Cache miss penalty}: $\sim$5--10\,ms per expert (page fault
    from mmap, dependent on OS page cache state and SSD latency).
  \item \textbf{Overhead amortisation}: with 10\% miss rate and
    $\sim$10 experts per token, approximately 1 expert per token incurs a
    cache miss, contributing 5--10\,ms of the total overhead.
\end{itemize}

The remaining overhead is attributable to the software indirection of cache
management versus direct tensor pointer access in the full-load scenario.

% ============================================================
% 6. LIMITATIONS
% ============================================================
\section{Limitations}
\label{sec:limitations}

\paragraph{Weight-level only.} The current ExpertFlow implementation operates
at the weight-loading level, providing a mechanism for efficiently accessing
expert tensors on demand. It does not yet integrate with an end-to-end
inference engine; the benchmarks in \Cref{sec:experiments:loading} measure
weight access patterns rather than full autoregressive generation with router
dispatch. Integration with frameworks such as mlx-lm is ongoing work.

\paragraph{Single-device scope.} ExpertFlow targets single-device inference on
consumer hardware. It does not address multi-device expert distribution, which
would enable even larger models through horizontal scaling. Expert
parallelism across multiple Apple Silicon devices connected via Thunderbolt is
a natural extension.

\paragraph{KV cache not addressed.} ExpertFlow optimises expert weight memory
but does not manage the key-value cache, which grows linearly with sequence
length and can itself become a memory bottleneck for long contexts. Combining
ExpertFlow with KV cache quantisation or paged
attention~\citep{kwon2023pagedattention} would provide a more complete memory
management solution.

\paragraph{Router dependency.} ExpertFlow's effectiveness depends on the
sparsity of expert activation. Models with low expert counts or high $k/N$
ratios would see diminished benefits. Similarly, adversarial inputs designed to
activate diverse expert sets could reduce cache hit rates, though such patterns
are uncommon in natural language.

\paragraph{Platform specificity.} While the core approach is
platform-agnostic, the performance characteristics we report are specific to
Apple Silicon's unified memory architecture. On discrete GPU systems, the
CPU$\rightarrow$GPU transfer cost for cache-missed experts would likely
increase the throughput penalty. Unified memory is essential to the
favourable tradeoff we observe.

% ============================================================
% 7. CONCLUSION & FUTURE WORK
% ============================================================
\section{Conclusion and Future Work}
\label{sec:conclusion}

We have presented ExpertFlow, a system that exploits the structural sparsity
of Mixture-of-Experts models to enable efficient inference on consumer
hardware. By memory-mapping model shards and dynamically loading experts
through an LRU cache, ExpertFlow reduces peak memory consumption by 93.5\%
(41.8\,GB $\rightarrow$ 2.7\,GB) for Qwen3-Next-80B-A3B with only a 12\%
throughput penalty. More significantly, ExpertFlow opens a path to serving
models that exceed device memory entirely---up to 2.75$\times$ available RAM
in our ongoing evaluations.

The 90\% cache hit rate observed across cache sizes confirms that MoE expert
activations exhibit strong temporal locality, a property that can be further
exploited through predictive strategies.

\subsection{Future Work}

\paragraph{mlx-lm integration.} Integrating ExpertFlow as a custom weight
provider in mlx-lm would enable full end-to-end generation with dynamic expert
loading, bridging the gap between our weight-level benchmarks and practical
inference.

\paragraph{Predictive expert pre-fetching.} Router weights are available
before expert dispatch. By running the router one layer ahead and
speculatively pre-loading predicted experts, cache miss latency can be hidden
behind computation. This is particularly effective for autoregressive
generation where subsequent tokens are predictable.

\paragraph{Correlated expert prefetching.} Analysis of expert co-activation
patterns may reveal stable clusters of experts that are frequently selected
together. Pre-fetching entire clusters on activation of any member would
reduce miss rates without requiring lookahead routing.

\paragraph{Multi-device expert sharding.} Distributing experts across multiple
Apple Silicon devices connected via Thunderbolt (up to 100\,Gbps) would enable
inference on arbitrarily large MoE models. Each device would host a shard of
the expert pool, with a lightweight coordination protocol for routing expert
requests.

\paragraph{Adaptive cache sizing.} Dynamically adjusting the LRU cache size
based on observed working set size and available memory pressure would allow
ExpertFlow to maximise throughput while respecting system constraints.

% ============================================================
% REFERENCES
% ============================================================
\bibliographystyle{plainnat}

\begin{thebibliography}{30}

\bibitem[Aminabadi et~al.(2022)]{aminabadi2022deepspeed}
Aminabadi, R.~Y., Rajbhandari, S., Awan, A.~A., Li, C., Li, D., Zheng, E.,
  Ruwase, O., Smith, S., Zhang, M., Rasley, J., and He, Y.
\newblock {DeepSpeed-Inference}: Enabling efficient inference of transformer
  models at unprecedented scale.
\newblock In \emph{SC '22}, 2022.

\bibitem[Apple(2024)]{apple2024m4}
Apple.
\newblock {Apple M4 family of chips}.
\newblock Technical report, Apple Inc., 2024.

\bibitem[Chowdhery et~al.(2022)]{chowdhery2022palm}
Chowdhery, A., Narang, S., Devlin, J., Bosma, M., et~al.
\newblock {PaLM}: Scaling language modeling with pathways.
\newblock \emph{arXiv preprint arXiv:2204.02311}, 2022.

\bibitem[{DeepSeek-AI}(2024)]{deepseekv2}
{DeepSeek-AI}.
\newblock {DeepSeek-V2}: A strong, economical, and efficient
  mixture-of-experts language model.
\newblock \emph{arXiv preprint arXiv:2405.04434}, 2024.

\bibitem[Dettmers et~al.(2022)]{dettmers2022gptint8}
Dettmers, T., Lewis, M., Belkada, Y., and Zettlemoyer, L.
\newblock {GPT3.int8()}: 8-bit matrix multiplication for transformers at scale.
\newblock In \emph{NeurIPS}, 2022.

\bibitem[Dettmers et~al.(2023)]{dettmers2023qlora}
Dettmers, T., Pagnoni, A., Holtzman, A., and Zettlemoyer, L.
\newblock {QLoRA}: Efficient finetuning of quantized language models.
\newblock In \emph{NeurIPS}, 2023.

\bibitem[Fedus et~al.(2022)]{fedus2022switch}
Fedus, W., Zoph, B., and Shazeer, N.
\newblock Switch Transformers: Scaling to trillion parameter models with simple
  and efficient sparsity.
\newblock \emph{JMLR}, 23(120):1--39, 2022.

\bibitem[Frantar \& Alistarh(2023)]{frantar2023sparsegpt}
Frantar, E. and Alistarh, D.
\newblock {SparseGPT}: Massive language models can be accurately pruned in one
  shot.
\newblock In \emph{ICML}, 2023.

\bibitem[Frantar et~al.(2023)]{frantar2023gptq}
Frantar, E., Ashkboos, S., Hoefler, T., and Alistarh, D.
\newblock {GPTQ}: Accurate post-training quantization for generative
  pre-trained transformers.
\newblock In \emph{ICLR}, 2023.

\bibitem[Gerganov(2023)]{gerganov2023llamacpp}
Gerganov, G.
\newblock {llama.cpp}.
\newblock \url{https://github.com/ggerganov/llama.cpp}, 2023.

\bibitem[Gu et~al.(2024)]{gu2024minillm}
Gu, Y., Dong, L., Wei, F., and Huang, M.
\newblock {MiniLLM}: Knowledge distillation of large language models.
\newblock In \emph{ICLR}, 2024.

\bibitem[Hinton et~al.(2015)]{hinton2015distilling}
Hinton, G., Vinyals, O., and Dean, J.
\newblock Distilling the knowledge in a neural network.
\newblock \emph{arXiv preprint arXiv:1503.02531}, 2015.

\bibitem[Hwang et~al.(2023)]{hwang2023semoe}
Hwang, C., Cui, W., Xiong, Y., Yang, Z., et~al.
\newblock Tutel: Adaptive mixture-of-experts at scale.
\newblock In \emph{MLSys}, 2023.

\bibitem[Hwang et~al.(2024)]{hwang2024pregated}
Hwang, R., Zhu, J., Chen, Y., et~al.
\newblock Pre-gated {MoE}: An algorithm-system co-design for fast and scalable
  mixture-of-experts inference.
\newblock \emph{arXiv preprint arXiv:2308.12066}, 2024.

\bibitem[Jiang et~al.(2024)]{jiang2024mixtral}
Jiang, A.~Q., Sablayrolles, A., Roux, A., et~al.
\newblock {Mixtral of Experts}.
\newblock \emph{arXiv preprint arXiv:2401.04088}, 2024.

\bibitem[Kwon et~al.(2023)]{kwon2023pagedattention}
Kwon, W., Li, Z., Zhuang, S., et~al.
\newblock Efficient memory management for large language model serving with
  {PagedAttention}.
\newblock In \emph{SOSP}, 2023.

\bibitem[Lepikhin et~al.(2021)]{lepikhin2021gshard}
Lepikhin, D., Lee, H., Xu, Y., Chen, D., Firat, O., Huang, Y., Krikun, M.,
  Shazeer, N., and Chen, Z.
\newblock {GShard}: Scaling giant models with conditional computation and
  automatic sharding.
\newblock In \emph{ICLR}, 2021.

\bibitem[Lin et~al.(2024)]{lin2024awq}
Lin, J., Tang, J., Tang, H., Yang, S., Chen, W.-M., Wang, W.-C., Xiao, G.,
  Dang, S., Gan, C., and Han, S.
\newblock {AWQ}: Activation-aware weight quantization for efficient {LLM}
  inference and deployment.
\newblock In \emph{MLSys}, 2024.

\bibitem[{LM Studio}(2024)]{lmstudio}
{LM Studio}.
\newblock {LM Studio}: Discover, download, and run local {LLMs}.
\newblock \url{https://lmstudio.ai}, 2024.

\bibitem[{MLX Contributors}(2023)]{mlx2023}
{MLX Contributors}.
\newblock {MLX}: An array framework for Apple Silicon.
\newblock \url{https://github.com/ml-explore/mlx}, 2023.

\bibitem[Rajbhandari et~al.(2022)]{rajbhandari2022deepspeedmoe}
Rajbhandari, S., Li, C., Yao, Z., Zhang, M., Aminabadi, R.~Y., Awan, A.~A.,
  Rasley, J., and He, Y.
\newblock {DeepSpeed-MoE}: Advancing mixture-of-experts inference and training
  to power next-generation {AI} scale.
\newblock In \emph{ICML}, 2022.

\bibitem[Shazeer et~al.(2017)]{shazeer2017outrageously}
Shazeer, N., Mirhoseini, A., Maciaszek, K., Davis, A., Le, Q., Hinton, G.,
  and Dean, J.
\newblock Outrageously large neural networks: The sparsely-gated
  mixture-of-experts layer.
\newblock In \emph{ICLR}, 2017.

\bibitem[Sheng et~al.(2023)]{sheng2023flexgen}
Sheng, Y., Zheng, L., Yuan, B., et~al.
\newblock {FlexGen}: High-throughput generative inference of large language
  models with a single {GPU}.
\newblock In \emph{ICML}, 2023.

\bibitem[Song et~al.(2023)]{song2023powerinfer}
Song, Y., Mi, Z., Xie, H., and Chen, H.
\newblock {PowerInfer}: Fast large language model serving with a
  consumer-grade {GPU}.
\newblock \emph{arXiv preprint arXiv:2312.12456}, 2023.

\bibitem[Sun et~al.(2023)]{sun2023wanda}
Sun, M., Liu, Z., Bair, A., and Kolter, J.~Z.
\newblock A simple and effective pruning approach for large language models.
\newblock \emph{arXiv preprint arXiv:2306.11695}, 2023.

\bibitem[Touvron et~al.(2023)]{touvron2023llama2}
Touvron, H., Martin, L., Stone, K., et~al.
\newblock {LLaMA~2}: Open foundation and fine-tuned chat models.
\newblock \emph{arXiv preprint arXiv:2307.09288}, 2023.

\bibitem[Wolf et~al.(2020)]{wolf2020transformers}
Wolf, T., Debut, L., Sanh, V., et~al.
\newblock Transformers: State-of-the-art natural language processing.
\newblock In \emph{EMNLP: System Demonstrations}, 2020.

\bibitem[Xue et~al.(2024)]{xue2024moeinf}
Xue, L., You, J., Zhang, Y., et~al.
\newblock {MoE-Infinity}: Activation-aware expert offloading for efficient
  {MoE} inference.
\newblock \emph{arXiv preprint arXiv:2401.14361}, 2024.

\bibitem[Zoph et~al.(2022)]{zoph2022stmoe}
Zoph, B., Bello, I., Kumar, S., Du, N., Huang, Y., Dean, J., Shazeer, N.,
  and Fedus, W.
\newblock {ST-MoE}: Designing stable and transferable sparse expert models.
\newblock \emph{arXiv preprint arXiv:2202.08906}, 2022.

\end{thebibliography}

% ============================================================
\appendix
\section{Reproducibility}
\label{app:reproducibility}

ExpertFlow is available at \url{https://github.com/jhammant/expertflow}
under the MIT licence. Benchmarks were run on macOS~26.3.1 with an Apple
M5~Max (128\,GB unified memory, 4\,TB NVMe SSD). Model weights were obtained
from Hugging Face Hub in safetensor format (4-bit quantised variants where
noted). The benchmarking scripts are included in the repository under
\texttt{benchmarks/}.

\end{document}
